% Italian -- what the reading editions' preamble leaves to the language.
% Input by lib/frank-preamble.tex as \input{\FrankLib/lang/\BookLang.tex};
% docs/languages.md section 6 lists the hooks every file here must define.
% The Latin-script model: no font of its own, no marks to strip, two passes.

% ---- the babel language and the switches ----------------------------------
\newcommand{\FrankLangName}{italian}
\FrankRTLfalse           % chunk left, gloss right; \pw needs no bidi isolate
\FrankHasBarefalse       % there is nothing to strip, so pass 3 would repeat pass 1
\FrankHasAltfalse        % no alternate face
\FrankHasReadingfalse
\FrankHasVertfalse

% ---- fonts ------------------------------------------------------------------
% The registry's tex.main is null for Italian: the text is set in the main
% roman the core preamble already loaded (TeX Gyre Pagella), the same face as
% the glosses, so an Italian chunk and its English gloss differ by size and
% colour only.  import=it brings babel's Italian hyphenation and quoting rules.
\babelprovide[import=it]{italian}

% ---- the vocabulary line -----------------------------------------------------
% \vb{infinitive}{sound}{first person present}{sound}{past participle}{sound}{meaning}
% A sound is the form with its stressed vowel marked, and only where the
% stress is not on the next-to-last syllable: \vb{prendere}{prèndere}{prendo}{}...
% (the open or closed e and o are the pronunciation line's, not the \vb's).
% A verb of the third person only gives its third (accadere: accade, p.r.
% accadde), as the dictionary does.
% The first person present is the one form that shows what the infinitive
% hides -- the -isc- of finisco, the stem of vado, esco, vengo, devo -- and the
% participle is what every compound past is made of (è andata, ha parlato),
% which is why the class label the gloss once carried (-are, -ire, -isc-) is
% gone: those two forms and the infinitive's own ending already say it.
% What they cannot say goes in ONE parenthesis after the meaning, items
% parted by "; ", and only when it is not the ordinary case:
%   aux. \pw{essere}           the perfect is made with essere (avere is the
%                              default and is never printed);
%   aux. \pw{avere}/\pw{essere} it takes both (finire, correre, cominciare);
%   p.r. \pw{venni}            the passato remoto, first person, and only when
%                              it is irregular -- the narrative past of every
%                              book here, and fece, disse, prese, venne cannot
%                              be walked back to their infinitives.
%   \vb{venire}{}{vengo}{}{venuto}{}{to come (aux. \pw{essere}; p.r. \pw{venni})}
% The reader takes the same pair of labels from the registry
% (lib/languages.json, "vb_labels": ["pres.", "p.p."], the forms in words as
% "vb_forms"); the two must agree, or the PDF and the reader of one book label
% the same form differently.  lib/newlang.py --check compares them.
\newcommand{\FrankVbPres}{pres.}
\newcommand{\FrankVbPast}{p.p.}

% ---- the two Lua filters ---------------------------------------------------
% Both are the identity: Italian carries no marks the bare pass would drop
% (the registry's `strip` is null) and its labels are already in Latin digits.
% They still have to exist, because the core calls them for every language.
\directlua{
  function frank_strip(s) return s end
  function frank_latin(s) return s end
}

% ---- the "How to read this" page -----------------------------------------------
\newcommand{\FrankHowTo}{%
Every passage comes twice.

\smallskip
\textbf{First} the whole sentence, and nothing else. Try to read it. Get as far
as you can and do not worry about what you cannot get — the point is the
attempt, made before any help arrives.

\smallskip
\textbf{Second} the same sentence in small chunks, Italian on the left, and on
the right three lines: how it is pronounced, then how to look each word up,
then what it means. Now you find out how much you had right. Italian is
written as it is spoken, so there is no third pass: the sentence you attempted
is already the sentence as it is printed.

\smallskip
The pronunciation line is given only where the spelling misleads, and is
empty otherwise: the stressed vowel carries an accent when it is not the
penultimate, open \textit{è}, \textit{ò} against closed \textit{é},
\textit{ó}; a dot marks the voiced \textit{ṡ} and \textit{ż}
(\textit{càṡa}, \textit{żèro}); \textit{gli} is written \textit{ʎ},
\textit{gn} is \textit{ɲ} and \textit{sc} before \textit{e} or \textit{i}
is \textit{ʃ}; a doubled consonant is held. Verbs are given as infinitive,
first person present, past participle, meaning — the past participle is the
one you cannot always guess — and a form carries an accent beside it only
where its stress is not on the next-to-last syllable (\textit{prèndere},
\textit{àbito}). What those three cannot tell you stands in brackets after
the meaning: \textit{aux. essere} where the perfect is made with
\textit{essere} (\textit{è andata}), \textit{aux. avere/essere} where it can
be made with either, and nothing at all where it takes only \textit{avere};
and the past historic, \textit{p.r.}, wherever it is irregular —
\textit{venni}, \textit{feci}, \textit{presi} — since that is the tense a
story is told in.
}
